\chapter{Trigeminal Neuralgia}

\section*{Definition}
Trigeminal neuralgia (TN) is a chronic pain condition characterized by recurrent, unilateral, brief, electric-shock-like paroxysmal pain in one or more divisions of the trigeminal nerve (CN V). The pain is triggered by innocuous stimuli (light touch, chewing, talking, brushing teeth) and is classified as classical (neurovascular compression), secondary (structural lesion), or idiopathic.

\section*{Pathophysiology}
Classical TN is most commonly caused by vascular compression of the trigeminal nerve root entry zone (REZ) in the prepontine cistern, typically by a tortuous superior cerebellar artery or anterior inferior cerebellar artery. Chronic pulsatile compression causes focal demyelination of large-diameter A-beta fibers at the REZ. Demyelinated axons become hyperexcitable and generate spontaneous ectopic impulses. Ephaptic transmission (cross-talk) between A-beta touch fibers and A-delta pain fibers allows light touch to trigger pain signals. Peripheral and central sensitization maintain the pain circuit. Secondary TN results from structural lesions such as vestibular schwannoma, meningioma, or multiple sclerosis plaques involving the trigeminal nerve.

\section*{Etiology}
\begin{itemize}
  \item \textbf{Classical (neurovascular):} Superior cerebellar artery compression, anterior inferior cerebellar artery, vertebral artery, pontine veins
  \item \textbf{Secondary (structural):} Vestibular schwannoma, meningioma, epidermoid cyst, arteriovenous malformation, trigeminal nerve schwannoma
  \item \textbf{Demyelinating:} Multiple sclerosis (bilateral TN is highly suggestive of MS)
  \item \textbf{Idiopathic:} No identifiable structural or vascular cause
  \item \textbf{Rare:} Basilar artery dolichoectasia, Paget disease with skull base involvement, connective tissue disease
\end{itemize}

\section*{Indications for Evaluation}
Seek care if:
\begin{itemize}
  \item Recurrent, brief, unilateral facial pain triggered by light touch or movement
  \item Pain refractory to first-line medications (carbamazepine, oxcarbazepine)
  \item Progressive sensory loss, facial numbness, or multiple divisions involved
  \item Bilateral facial pain or age < 40 years (raise suspicion for multiple sclerosis)
\end{itemize}

\section*{Related Symptoms}
\begin{itemize}
  \item Brief (seconds to 2 minutes), lancinating electric-shock pain in V2 (maxillary) or V3 (mandibular) divisions most commonly
  \item Unilateral pain triggered by non-noxious stimuli (touch, wind, shaving, eating, talking)
  \item Refractory period after a paroxysm
  \item No sensory loss or motor deficit in classical TN
  \item Autonomic symptoms (lacrimation, rhinorrhea) may occur but are less prominent than in SUNHA
  \item Tic douloureux: facial grimacing during a pain paroxysm
\end{itemize}
\textbf{Note:} Sensory loss or hypesthesia in the trigeminal distribution suggests secondary TN (structural lesion or MS) and warrants MRI with contrast.

\section*{Self-Care / Home Management}
\begin{itemize}
  \item Avoid triggers: gentle face washing, lukewarm food, soft diet, electric razor
  \item Eat on the unaffected side when possible
  \item Room-temperature food and drink (avoid extremes)
  \item Stress reduction to decrease paroxysm frequency
  \item Medical alert bracelet if on carbamazepine or oxcarbazepine (risk of hyponatremia)
\end{itemize}

\section*{Diagnostic Approach}
\begin{itemize}
  \item Clinical diagnosis using the International Classification of Headache Disorders (ICHD-3) criteria
  \item MRI brain with high-resolution CISS or FIESTA sequences to visualize neurovascular compression
  \item MRI with contrast for secondary causes: vestibular schwannoma, meningioma, demyelinating plaques (MS protocol)
  \item MR angiography/venography to assess vascular anatomy
  \item Trigeminal reflex testing (blink reflex, masseter inhibitory reflex) in atypical cases
  \item Lumbar puncture with oligoclonal band analysis if multiple sclerosis is suspected
\end{itemize}

\section*{Treatment / Management Overview}
\begin{itemize}
  \item First-line pharmacotherapy: Carbamazepine (200--1200 mg/day) or oxcarbazepine (600--1800 mg/day)
  \item Second-line: Lamotrigine, topiramate, gabapentin, pregabalin, baclofen
  \item Acute rescue: Lidocaine or fosphenytoin IV for severe paroxysmal attacks
  \item Microvascular decompression (MVD): Posterior fossa surgery to separate the compressing vessel from the trigeminal nerve; highest long-term success rate
  \item Gamma Knife radiosurgery: Stereotactic radiation to the trigeminal nerve REZ
  \item Percutaneous procedures: Balloon compression, glycerol rhizotomy, radiofrequency thermocoagulation
  \item Treatment of underlying cause (MS disease-modifying therapy, tumor resection)
\end{itemize}

\clearpage
